\section{Related Work}
\label{sec:related-work}


{\bf Offline Policy Distillation:} Most closely related to our work are the recent advances in learning policies from offline environment interaction data with transformers, which we have been referring to as Policy Distillation (PD). Initial PD architectures such as Decision Transformer~\citep{chen2021dt} and Trajectory Transformer~\citep{janner2021tto} showed that transformers can learn single-task policies from offline data. Subsequently the Multi-Game Decision Transformer (MGDT)~\citep{lee2022mgdt} and Gato~\citep{reed2022gato} showed that PD architectures can also learn multi-task same domain and cross-domain policies, respectively. Importantly, these prior methods use contexts substantially smaller than an episode length, which is likely the reason in-context RL was not observed in these works. Instead, they rely on alternate ways to adapt to new tasks - MGDT finetunes the model parameters while Gato gets prompted with expert demonstrations to adapt to downstream tasks. \name adapts in-context without finetuning and does not rely on demonstrations.
Finally, a number of recent works have explored more generalized PD architectures~\citep{furuta2022gdt}, prompt conditioning~\citep{xu2022pdt}, and online gradient-based RL~\citep{zheng2022odt}.
\looseness=-1


{\bf Meta Reinforcement Learning:} \name falls into the category of methods that learn to reinforcement learn, also known as meta-RL. Specifically, \name is an in-context offline meta-RL method. This general idea of learning the policy improvement process has a long history in reinforcement learning, but has been limited to meta-learning hyper-parameters until recently~\citep{ishii2002control}. In-context deep meta-RL methods introduced by~\citet{wang2016metarl} and~\citet{duan2016rl2} are usually trained in the online setting by maximizing multi-episodic value functions with memory-based architectures through environment interactions. Another common approach to online meta-RL includes optimization-based methods that find good network parameter initializations for meta-RL~\citep{hochreiter2001learning, finn17maml, nichol2018reptile} and adapt by taking additional gradient steps. Like other in-context meta-RL approaches, \name is gradient-free - it adapts to downstream tasks without updating its network parameters. Recent works have proposed learning to reinforcement learn from offline datasets, or offline meta-RL, using Bayesian RL~\citep{dorfman21omrl} and optimization-based meta-RL~\citep{mitchell21macaw}. Given the difficulty of offline meta-RL, \citet{pong2022offline} proposed a hybrid offline-online strategy for meta-RL. 

{\bf In-Context Learning with Transformers:} In this work, we make the distinction between in-context learning and incremental or in-context learning. In-context learning involves learning from a provided prompt or demonstration while incremental in-context learning involves learning from one's own behavior through trial and error.
While many recent works have demonstrated the former, it is much less common to see methods that exhibit the latter. Arguably, the most impressive demonstrations of in-context learning to date have been shown in the text completion setting~\citep{radford2018improving, chen2020igpt, brown2020gpt3} through prompt conditioning. Similar methodology was recently extended to show powerful composability properties in text-conditioned image generation~\citep{yu2022parti}. Recent work showed that transformers can also learn simple algorithm classes, such as linear regression, in-context in a small-scale setting~\citep{garg2022incontext}. Like prior in-context learning methods,~\citet{garg2022incontext} required initializing the transformer prompt with expert examples. While the aforementioned approaches were examples of in-context learning, a recent work~\citep{chen2022optformer} demonstrated incremental in-context learning for hyperparameter optimization by treating hyperparameter optimization as a sequential prediction problem with a score function.